\pagestyle{appendix}
\appendix
\renewcommand{\thesubsection}{\Alph{subsection}}

\subsection{Complete Source File Directory}
\label{app:source-files}

This appendix lists every source file from national-platform/ that was used to
build this paper, organized by directory. For each file, the filename, line
count, and a description of its content are provided.

\noindent\textbf{national-platform/21cfr312\_adapt/}
\begin{itemize}[nosep]
  \item 01\_preamble\_scope\_definitions.tex (427 lines) - Subpart A: General
    Provisions including 22 Physical AI definitions
  \item 02\_ind\_content\_phases.tex (274 lines) - Subpart B: IND Application
    including Phase 0 simulation validation
  \item 03\_protocol\_amendments\_reporting.tex (328 lines) - Information
    amendments, safety reporting, Subpart C start
  \item 04\_annual\_reports\_withdrawal.tex (463 lines) - Inactive status,
    meetings, Subpart D: Sponsor/Investigator responsibilities
  \item 05\_clinical\_holds\_appendices\_closing.tex (783 lines) - Subparts
    E through J including new Physical AI Systems subpart
\end{itemize}

\noindent\textbf{national-platform/21cfr50\_adapt/}
\begin{itemize}[nosep]
  \item 01\_preamble\_scope\_definitions\_consent.tex (438 lines) - Subparts
    A and B: Scope, definitions, informed consent
  \item 02\_irb\_review\_pediatric.tex (236 lines) - Subparts C and D:
    Physical AI safety requirements, pediatric protections
  \item 03\_additional\_safeguards\_closing.tex (73 lines) - Glossary and
    document closing
\end{itemize}

\noindent\textbf{national-platform/ich\_e6r3\_adapt/}
\begin{itemize}[nosep]
  \item 01\_preamble\_principles\_investigator.tex (337 lines) - Sections 1-2:
    Principles and investigator responsibilities
  \item 02\_sponsor\_responsibilities.tex (346 lines) - Section 3: Sponsor
    responsibilities in Physical AI trials
  \item 03\_data\_governance.tex (228 lines) - Section 4: Data governance
    for Physical AI trials
  \item 04\_appendices\_glossary.tex (389 lines) - Appendices A-C and glossary
\end{itemize}

\noindent\textbf{national-platform/new\_trial\_psl/}
\begin{itemize}[nosep]
  \item 01\_preamble\_and\_sb1042.tex (468 lines) - SB 1042: Trial
    Authorization and Site Establishment Act
  \item 02\_ab2847\_patient\_rights.tex (326 lines) - AB 2847: Patient Rights
    and Robotic Safety Act
  \item 03\_sb892\_data\_protection.tex (307 lines) - SB 892: Data Protection
    and Transparency Act
  \item 04\_sf\_municipal\_code.tex (262 lines) - San Francisco Municipal Code
    Update
  \item 05\_title22\_regulations.tex (322 lines) - Title 22 Regulations
  \item 06\_fda\_compliance\_guide.tex (317 lines) - FDA National Compliance
    Guide
  \item 07\_building\_code.tex (287 lines) - Building Code Standards
  \item 08\_premises\_code.tex (276 lines) - Premises Code
  \item 09\_parking\_transportation.tex (250 lines) - Parking and
    Transportation Standards
  \item 10\_activation\_sops.tex (272 lines) - Activation and Standard
    Operating Procedures
  \item 11\_emergency\_preparedness.tex (290 lines) - Emergency Preparedness
    Plan
\end{itemize}

\noindent\textbf{national-platform/patient\_journey/}
\begin{itemize}[nosep]
  \item 01\_preamble\_introduction\_stages1to4.tex (296 lines) - Introduction
    and pre-trial stages
  \item 02\_stages5to8.tex (358 lines) - Treatment stages and outcomes
  \item 03\_stages9to10\_discussion\_closing.tex (222 lines) - Closeout,
    discussion, and conclusions
\end{itemize}

\noindent\textbf{national-platform/patient\_robot/}
\begin{itemize}[nosep]
  \item 01\_preamble\_robots1to5.tex (197 lines) - Patient instructions for
    robot types 1-5
  \item 02\_robots6to10\_closing.tex (173 lines) - Patient instructions for
    robot types 6-10
\end{itemize}

\noindent\textbf{national-platform/national\_mcp/}
\begin{itemize}[nosep]
  \item chunk1\_preamble\_and\_introduction.tex (181 lines) - Introduction
    and fragmentation analysis
  \item chunk2\_methods\_and\_results\_part1.tex (284 lines) - Methods and
    five-server architecture
  \item chunk3\_results\_part2.tex (224 lines) - Safety modules and
    integration
  \item chunk4\_discussion\_and\_conclusion.tex (322 lines) - National
    deployment and governance
\end{itemize}

\noindent\textbf{national-platform/federated\_learning/}
\begin{itemize}[nosep]
  \item main\_chunk1\_preamble\_intro\_methods.tex (215 lines) - Introduction
    and methods
  \item main\_chunk2\_results\_architecture\_pillars\_examples.tex
    (318 lines) - Architecture and five pillars
  \item main\_chunk3\_results\_peerreview\_trust\_analytics.tex
    (243 lines) - Peer review and analytics
  \item main\_chunk4\_discussion\_limitations\_conclusion.tex (154 lines) -
    Discussion and conclusions
\end{itemize}

\noindent\textbf{national-platform/usl\_standard/}
\begin{itemize}[nosep]
  \item 01\_preamble\_introduction\_framework.tex (249 lines) - Introduction
    and USL methodology
  \item 02\_results\_discussion\_closing.tex (227 lines) - Results and
    discussion
\end{itemize}

\noindent\textbf{national-platform/research\_a/}
\begin{itemize}[nosep]
  \item 01\_research\_a\_part1.txt (113 lines) - U.S. government framework
    and AI clinical trial regulation
  \item 02\_research\_a\_part2.txt (166 lines) - Reference URLs and citations
\end{itemize}

\noindent\textbf{national-platform/research\_b/}
\begin{itemize}[nosep]
  \item 01\_research\_b\_part1.txt (95 lines) - California and federal
    regulatory analysis
  \item 02\_research\_b\_part2.txt (109 lines) - Reference URLs and citations
\end{itemize}

\newpage

\subsection{Glossary of Physical AI Terms}
\label{app:glossary}

This glossary compiles unified Physical AI terms used throughout the paper,
drawn from the glossaries in all three adapted regulatory standards. Terms are
listed alphabetically with definitions and primary source references.

\begin{longtable}{>{\raggedright\arraybackslash}p{4.5cm} >{\raggedright\arraybackslash}p{11cm}}
\toprule
\textbf{Term} & \textbf{Definition} \\
\midrule
\endfirsthead
\toprule
\textbf{Term} & \textbf{Definition} \\
\midrule
\endhead
\bottomrule
\endfoot
Agentic AI & AI systems capable of autonomous multi-step task execution, coordinating complex clinical workflows without continuous human direction. \\
Autonomous Mode & Operation of a Physical AI system where AI algorithms direct robotic actions without real-time human input, subject to safety boundaries and override capability. \\
Calibration & The process of adjusting Physical AI system parameters (position sensors, force sensors, actuators) to meet specified accuracy and precision standards. \\
Cobot & Collaborative robot designed to work safely alongside humans in shared workspaces, with force-limiting and collision-detection safety features. \\
Cybersecurity Incident & Any unauthorized access, disruption, or modification of Physical AI systems that may affect patient safety, data integrity, or trial operations. \\
Digital Twin & A computational model replicating a physical patient, tumor, or system for simulation, prediction, and treatment planning purposes. \\
Emergency Stop (E-Stop) & A safety mechanism that immediately halts all Physical AI system motion and operations, available through hardware buttons, software commands, and wireless triggers. \\
Federated Learning & A machine learning approach where models are trained across decentralized data sources without exchanging raw patient data, preserving privacy while enabling model improvement. \\
Humanoid Robot & A robot with human-like form factor designed for general-purpose tasks including patient guidance, logistics support, and facility operations. \\
MCP Server & A Model Context Protocol server providing standardized AI-to-tool communication for Physical AI system coordination across clinical trial operations. \\
Needle-Placement Robot & A robotic system that guides needles for biopsies and targeted drug delivery with precision exceeding manual capability. \\
Phase 0 Simulation & A mandatory pre-enrollment phase where Physical AI systems demonstrate safety and preliminary efficacy through validated simulations before human subjects contact. \\
Physical AI Adverse Event & Any untoward occurrence during a clinical investigation that is attributable to Physical AI system operation, including robot malfunctions, AI decision errors, and sensor failures. \\
Physical AI Operator & A qualified individual responsible for monitoring and controlling Physical AI systems during clinical procedures, with training per adapted ICH E6(R3) requirements. \\
Physical AI Safety Matrix & A pre-procedure verification checklist ensuring all Physical AI safety systems (mechanical, sensor, E-stop, communication, AI model) are operational before patient interaction. \\
Physical AI System & An integrated robotic and AI platform designed for direct or indirect patient interaction in clinical trial settings, encompassing hardware, software, and AI components. \\
PSL & Physical AI Standard Level, a three-dimension framework evaluating trial site compliance across legislative authorization, regulatory compliance, and operational readiness. \\
Robotic Procedure & Any clinical intervention performed wholly or partly by a Physical AI system, including surgical, diagnostic, therapeutic, and support activities. \\
Social Companion Robot & A robot designed to provide emotional support, information delivery, and well-being monitoring for patients during clinical trial participation. \\
Supervised Mode & Operation of a Physical AI system where human operators maintain continuous real-time control over all robotic actions. \\
Surgical Robot & A robotic system designed for surgical procedures, operated under direct surgeon control (teleoperated) or with supervised autonomous capabilities. \\
USL & Unification Standard Level, a four-dimension 1.0-10.0 scoring framework evaluating robot readiness for multi-site clinical trial deployment. \\
USL Band & Classification tier based on composite USL score: Exemplary (9.0-10.0), Advanced (7.0-8.9), Intermediate (5.0-6.9), Foundational (3.0-4.9), Initial (1.0-2.9). \\
\end{longtable}

\newpage

\subsection{Regulatory Cross-Reference Matrix}
\label{app:cross-reference}

This matrix shows how each section of this paper maps to source documents,
source files, applicable regulatory standards, and key references.

\begin{longtable}{>{\raggedright\arraybackslash}p{2.6cm} >{\raggedright\arraybackslash}p{3cm} >{\raggedright\arraybackslash}p{3.6cm} >{\raggedright\arraybackslash}p{2.3cm} >{\raggedright\arraybackslash}p{3.5cm}}
\toprule
\textbf{Section} & \textbf{Source Docs} & \textbf{Source Files} & \textbf{Standards} & \textbf{Key References} \\
\midrule
\endfirsthead
\toprule
\textbf{Section} & \textbf{Source Docs} & \textbf{Source Files} & \textbf{Standards} & \textbf{Key References} \\
\midrule
\endhead
\bottomrule
\endfoot
\S~1 Intro & All source docs & All directories & All three & \cite{tufts-delay, fda-clinical-trials-guidance, nci-clinical-trials} \\
\S~2 Gov Framework & Research A & research\_a/ (2 files) & CFR 312, 50, ICH & \cite{cfr-part-312, cfr-part-50, ich-e6r3, hipaa} \\
\S~3 Regulatory & Research B & research\_b/ (2 files) & All three & \cite{ich-adapt, cfr50-adapt, cfr312-adapt} \\
\S~4 ICH E6(R3) & ICH Adaptation & ich\_e6r3\_adapt/ (4 files) & ICH E6(R3) & \cite{ich-adapt, ich-e6r3, usl-standard} \\
\S~5 CFR Part 50 & CFR 50 Adapt & 21cfr50\_adapt/ (3 files) & 21 CFR 50 & \cite{cfr50-adapt, cfr-part-50} \\
\S~6 CFR Part 312 & CFR 312 Adapt & 21cfr312\_adapt/ (5 files) & 21 CFR 312 & \cite{cfr312-adapt, cfr-part-312} \\
\S~7 PSL/USL & USL Standard, Site Docs & usl\_standard/ (2), new\_trial\_psl/ (11) & All three & \cite{usl-standard, site-docs} \\
\S~8 Site Est. & Site Docs & new\_trial\_psl/ (11 files) & All three & \cite{site-docs, ich-adapt, cfr50-adapt} \\
\S~9 Journey & Patient Journey & patient\_journey/ (3 files) & All three & \cite{patient-journey-paper, usl-standard} \\
\S~10 Patient Inst & Patient Inst & patient\_robot/ (2 files) & CFR 50, ICH & \cite{patient-instructions-paper, cfr50-adapt} \\
\S~11 MCP & National MCP & national\_mcp/ (4 files) & All three & \cite{national-mcp-paper, mcp-protocol, mcp-repo} \\
\S~12 FL & FL Paper & federated\_learning/ (4 files) & HIPAA, ICH & \cite{fl-paper, fl-repo, federated-learning-healthcare} \\
\S~13 Financial & Journey, Site & patient\_journey/, new\_trial\_psl/ & N/A & \cite{tufts-delay, nci-trials-paying} \\
\S~14 Strategy & All sources & All directories & All three & \cite{main-repo, mcp-repo, fl-repo} \\
\S~15 Discussion & All sources & All directories & All three & All references \\
\S~16 Conclusion & All sources & All directories & All three & All references \\
\end{longtable}

\newpage

\subsection{PSL and USL Scoring Reference}
\label{app:scoring}

\noindent\textbf{PSL Scoring Reference}

\begin{table}[H]
\centering
\caption{PSL Three-Dimension Scoring Criteria}
\label{tab:psl-scoring}
\small
\begin{tabularx}{\textwidth}{c l X l}
\toprule
\textbf{Dim} & \textbf{Name} & \textbf{Key Criteria} & \textbf{Status} \\
\midrule
1 & Legislative & State authorization act, patient rights bill, data protection act & Pass/Fail \\
2 & Regulatory & Title 22 compliance, FDA guide, building/premises codes & Pass/Fail \\
3 & Operational & Activation SOPs, emergency plan, staff training, equipment verified & Pass/Fail \\
\bottomrule
\end{tabularx}
\end{table}

\noindent A site must achieve Passing status on all three dimensions before
Physical AI trial operations may begin.

\vspace{0.5cm}

\noindent\textbf{USL Scoring Reference}

\begin{table}[H]
\centering
\caption{USL Four-Dimension Scoring Summary}
\label{tab:usl-scoring}
\small
\begin{tabularx}{\textwidth}{c l c X}
\toprule
\textbf{Dim} & \textbf{Name} & \textbf{Weight} & \textbf{Scoring Factors} \\
\midrule
A & Simulation Switching & 25\% & Framework count, GPU support, transfer testing, model formats \\
B & AI Integration & 25\% & LLM planning, VLA support, diffusion policy, MCP, agentic AI \\
C & Cross-Robot Sharing & 25\% & ONNX export, ROS 2, checkpoints, inter-org transfer \\
D & Clinical Trial Collab & 25\% & FDA clearance, federated learning, audit trail, multi-site \\
\bottomrule
\end{tabularx}
\end{table}

\noindent\textbf{USL Band Classification:} Exemplary (9.0-10.0), Advanced
(7.0-8.9), Intermediate (5.0-6.9), Foundational (3.0-4.9), Initial (1.0-2.9).
Final score computed as weighted average of four dimension scores, clamped to
range 1.0-10.0.

\newpage

\subsection{Simulation Evidence Summary}
\label{app:simulation}

\begin{table}[H]
\centering
\caption{Single-Patient Journey Evidence Summary}
\label{tab:single-patient}
\small
\begin{tabularx}{\textwidth}{l X}
\toprule
\textbf{Parameter} & \textbf{Value} \\
\midrule
Patient diagnosis & Stage IIIB non-small cell lung cancer (NSCLC) \\
Trial duration & 1,120 days (prescreening through closeout) \\
Pipeline stages & 10 stages, fully autonomous management \\
Robot types involved & Surgical, cobots, imaging, companion, rehab exoskeleton \\
Regulatory coverage & All three adapted standards at every stage \\
ASCII diagrams & 30 three-perspective diagrams \\
Cost analysis & Per-stage cost savings documented \\
Source & \cite{patient-journey-paper} \\
\bottomrule
\end{tabularx}
\end{table}

\begin{table}[H]
\centering
\caption{24-Hour Multi-Patient Simulation Evidence Summary}
\label{tab:multi-patient}
\small
\begin{tabularx}{\textwidth}{l X}
\toprule
\textbf{Parameter} & \textbf{Value} \\
\midrule
Patients treated & 168 patients \\
Cancer types & 15 different cancer types \\
Robots deployed & 29 robots across multiple categories \\
System uptime & 99.7 percent continuous operation \\
Patient harm events & Zero patient harm events \\
Adverse events & 7 events documented per ICH E6(R3) \\
Operating duration & 24 continuous hours \\
Documentation & 72 ASCII diagrams, 11 regulatory documents \\
PSL validation & All three dimensions assessed \\
Source & \cite{site-docs} \\
\bottomrule
\end{tabularx}
\end{table}
